\documentclass[a4paper,12pt]{article}
\usepackage[T2A]{fontenc}
\usepackage[utf8]{inputenc}
\usepackage[russian]{babel}
\usepackage{amsmath,amssymb,mathtools}
\usepackage{helvet}
\renewcommand{\familydefault}{\sfdefault}
\usepackage{geometry}
\geometry{left=23mm,right=23mm,top=22mm,bottom=22mm}
\usepackage{microtype}
\usepackage{icomma}
\usepackage{tikz}
\usetikzlibrary{arrows.meta,positioning,shapes.geometric,calc}
\usepackage{booktabs,tabularx}
\usepackage{enumitem}
\usepackage{setspace}
\usepackage{fancyhdr}
\usepackage{xcolor}
\usepackage{hyperref}
\definecolor{accent}{HTML}{008080}
\definecolor{darkaccent}{HTML}{00585A}
\definecolor{textgray}{HTML}{35393B}
\definecolor{muted}{HTML}{677076}
\definecolor{warn}{HTML}{AA493C}
\hypersetup{colorlinks=true,linkcolor=darkaccent}
\setstretch{1.12}
\setlength{\parindent}{0pt}
\setlength{\parskip}{3pt}
\setcounter{secnumdepth}{1}
\setlist[itemize]{leftmargin=1.7em,itemsep=2pt,topsep=3pt}
\setlist[enumerate]{leftmargin=1.8em,itemsep=7pt,topsep=4pt}
\newcommand{\lead}[1]{\textcolor{darkaccent}{\textbf{#1}}}
\newcommand{\beware}[1]{\textcolor{warn}{\textbf{#1}}}
\newcommand{\sect}[1]{\section*{#1}\vspace{-3mm}}
\pagestyle{fancy}
\fancyhf{}
\fancyhead[L]{\small\textcolor{muted}{Основы алгебры}}
\fancyhead[R]{\small\textcolor{muted}{Володя Хачатурян}}
\fancyfoot[C]{\small\textcolor{muted}{\thepage}}
\renewcommand{\headrulewidth}{0pt}
\setlength{\headheight}{14pt}
\begin{document}
\thispagestyle{empty}
\begin{center}
\vspace*{24mm}
{\color{darkaccent}\Large МАТЕМАТИКА\par}
\vspace{23mm}
{\Huge\bfseries Чек-лист\par}
\vspace{15mm}
{\LARGE\color{accent} Пропорции и преобразование\\[3pt] алгебраических выражений\par}
\vspace{16mm}
{\large Базовые навыки для математики и физики\par}
\vfill
{\Large Володя Хачатурян\par}
\vspace{12mm}
{\large 26 октября 2026 г.\par}
\vspace{19mm}
{\normalsize Лёвин Артём Александрович\par}
{\small эксклюзивно для Володи Хачатуряна\par}
\vspace{18mm}
\begin{tikzpicture}
\draw[darkaccent,line width=0.85pt] (-5.5,0) .. controls (-4.1,0.8) and (-2.4,-0.8) .. (-0.9,0) .. controls (0.4,0.7) and (2,-0.7) .. (3.3,0) .. controls (4.2,0.45) and (4.9,0.1) .. (5.5,0);
\fill[darkaccent] (-0.9,0) circle (1.4pt);
\end{tikzpicture}
\vspace*{2mm}
\end{center}
\clearpage

\sect{1. Пропорция и её свойства}
\lead{Пропорция} --- равенство двух отношений:
\[
\frac{a}{b}=\frac{c}{d},\qquad b\ne0,\quad d\ne0.
\]

С пропорцией можно выполнять несколько полезных преобразований.
\begin{enumerate}[label=\arabic*.]
\item \textbf{Менять части местами:} $\displaystyle\frac{c}{d}=\frac{a}{b}$.
\item \textbf{Переворачивать обе дроби одновременно:}
$\displaystyle\frac{b}{a}=\frac{d}{c}$, если $a\ne0$ и $c\ne0$.
\item \textbf{Умножать крест-накрест:}
$\displaystyle ad=bc$.
\end{enumerate}
\lead{Диагональный перенос} --- удобное название для действий с множителями, основанных на умножении или делении \emph{обеих частей} равенства на одно и то же ненулевое число.
\[
\frac{a}{b}=\frac{c}{d}
\quad\Longrightarrow\quad
 a=\frac{bc}{d},\qquad
 c=\frac{ad}{b}.
\]
\textbf{Пример.} Из $\displaystyle\frac{x}{6}=\frac{5}{3}$ получаем
$\displaystyle x=\frac{6\cdot5}{3}=10$.

\lead{Единица в дроби.} Любое число $A$ можно записать как $\displaystyle\frac{A}{1}$. Единица в знаменателе обычно опускается; единицу в числителе убирать нельзя:
\[
\frac{A}{1}=A,\qquad \frac{1}{A}\ne A\quad\text{в общем случае}.
\]
\beware{Важно:} при сокращении числитель и знаменатель делятся на один и тот же ненулевой множитель; после сокращения вместо него может остаться единица.

\sect{2. Выражение величины из физической формулы}
Рассмотрим закон всемирного тяготения:
\[
F=G\frac{m_1m_2}{R^2}.
\]
Здесь $G>0$ --- гравитационная постоянная, $m_1,m_2>0$ --- массы, $R>0$ --- расстояние, $F>0$ --- модуль силы.

\textbf{Выразим первую массу.} Запишем левую часть как дробь и преобразуем:
\[
\frac{F}{1}=\frac{Gm_1m_2}{R^2}
\quad\Longrightarrow\quad FR^2=Gm_1m_2
\quad\Longrightarrow\quad
\boxed{m_1=\frac{FR^2}{Gm_2}}.
\]
Аналогично
\[
\boxed{m_2=\frac{FR^2}{Gm_1}},\qquad
\boxed{R^2=\frac{Gm_1m_2}{F}},\qquad
\boxed{R=\sqrt{\frac{Gm_1m_2}{F}}}.
\]
\beware{Не путайте:} $R^2=A\Rightarrow R=\sqrt A$ (при $R>0$); $2R=A\Rightarrow R=A/2$.

\sect{3. Четыре базовые схемы для нахождения неизвестного}
Вместо угадывания действий полезно определять, какую роль играет неизвестная буква.
\begin{center}
\renewcommand{\arraystretch}{1.7}
\begin{tabularx}{\linewidth}{@{}>{\raggedright\arraybackslash}p{32mm}X@{}}
\toprule
\textbf{Запись} & \textbf{Как найти неизвестное}\\
\midrule
$a+b=c$ & $a=c-b$,\quad $b=c-a$\\
$a-b=c$ & $a=b+c$,\quad $b=a-c$\\
$ab=c$ & $a=\dfrac cb$ при $b\ne0$;\quad $b=\dfrac ca$ при $a\ne0$\\
$\dfrac ab=c$ & $a=bc$;\quad $b=\dfrac ac$ при $c\ne0$ и $a\ne0$\\
\bottomrule
\end{tabularx}
\end{center}
В последней строке исходно требуется $b\ne0$. Если $c=0$, уравнение $a/b=0$ означает $a=0$ и значение $b$ из него однозначно найти нельзя.

\textbf{Пример с вычитаемым.} Если $17-b=9$, то $b=17-9=8$.
Проверка: $17-8=9$.

\sect{4. Линейное уравнение: переносы и знаки}
\textbf{Пример с занятия:}
\[
\begin{aligned}
4x+1&=10x-3,\\
4x-10x&=-3-1,\\
-6x&=-4,\\
x&=\frac{-4}{-6}=\boxed{\frac23}.
\end{aligned}
\]
\lead{Что помнить:}
\begin{itemize}
\item При переносе слагаемого через знак равенства его знак меняется.
\item В произведении $-6x$ коэффициент равен $-6$: знак относится к числу.
\item Для нахождения $x$ разделите \emph{обе части} на коэффициент при $x$.
\end{itemize}
\textbf{Проверка:} $4\cdot\frac23+1=\frac{11}{3}$ и $10\cdot\frac23-3=\frac{11}{3}$.

\clearpage

\begin{center}
\vspace*{5mm}
{\LARGE\color{darkaccent}\bfseries Алгоритм: дробное уравнение}\par
\vspace{9mm}
\begin{tikzpicture}[
>=Latex,
flow/.style={rounded corners=2pt,draw=darkaccent,line width=.7pt,align=center,text width=103mm,minimum height=14mm,inner sep=4pt,font=\normalsize},
dec/.style={diamond,aspect=2.15,draw=darkaccent,line width=.7pt,align=center,inner sep=2.8pt,text width=42mm,font=\small},
arr/.style={->,line width=.8pt,draw=darkaccent}]
\node[flow] (start) {Определите знаменатель и запишите ОДЗ: он не равен нулю};
\node[flow,below=10mm of start] (clear) {Уберите знаменатели равносильным преобразованием на ОДЗ};
\node[flow,below=10mm of clear] (solve) {Раскройте скобки, соберите слагаемые и найдите $x$};
\node[dec,below=14mm of solve] (check) {Корень удовлетворяет ОДЗ и исходному уравнению?};
\node[flow,below left=16mm and 2mm of check,text width=47mm,minimum height=16mm] (reject) {Нет: исключите найденное значение};
\node[flow,below right=16mm and 2mm of check,text width=47mm,minimum height=16mm] (accept) {Да: запишите ответ};
\draw[arr] (start)--(clear);
\draw[arr] (clear)--(solve);
\draw[arr] (solve)--(check);
\draw[arr] (check.south west)--node[above left,font=\small]{нет}(reject.north);
\draw[arr] (check.south east)--node[above right,font=\small]{да}(accept.north);
\end{tikzpicture}
\end{center}
\clearpage

\sect{5. Уравнения с неизвестным в числителе}
Если $\displaystyle\frac{4x+1}{3}=10$, умножим обе части на $3$:
\[
4x+1=30\quad\Longrightarrow\quad4x=29\quad\Longrightarrow\quad
\boxed{x=\frac{29}{4}}.
\]
\sect{6. Неизвестное в знаменателе: три способа}
Решим пример, разобранный на доске:
\[
\frac{3}{4x+1}=10,\qquad 4x+1\ne0\ \left(x\ne-\frac14\right).
\]
\textbf{Способ 1. Перевернуть обе части, затем выразить числитель:}
\[
\frac{4x+1}{3}=\frac1{10}\Longrightarrow 4x+1=\frac3{10}.
\]
\textbf{Способ 2. Сразу найти знаменатель по формуле $a/b=c$:}
\[
4x+1=\frac{3}{10}.
\]
\textbf{Способ 3. Перемножить крест-накрест:}
\[
\frac3{4x+1}=\frac{10}{1}
\Longrightarrow 3=10(4x+1)
\Longrightarrow 3=40x+10.
\]
Во всех случаях дальше получаем:
\[
4x=\frac3{10}-1=-\frac7{10},\qquad
\boxed{x=-\frac7{40}}.
\]
\lead{Проверка:}
\[
4\left(-\frac7{40}\right)+1=\frac3{10}\ne0,
\qquad \frac{3}{3/10}=10.
\]
\beware{Обязательно сохраняйте минус:} $\frac3{10}-1=-\frac7{10}$.

\sect{7. Отрицательный коэффициент и знак дроби}
\textbf{Другой пример с занятия:}
\[
\frac{8}{-2x+3}=2,\qquad -2x+3\ne0.
\]
После умножения обеих частей на знаменатель:
\[
8=2(-2x+3)=-4x+6
\Longrightarrow 2=-4x
\Longrightarrow \boxed{x=-\frac12}.
\]
\textbf{Проверка:} $-2\cdot(-\frac12)+3=4$ и $\frac84=2$.

Знак минус можно перенести из знаменателя перед всю дробь или в числитель:
\[
\frac{a}{-b}=-\frac ab=\frac{-a}{b},\qquad b\ne0.
\]
Для составного числителя скобки сохраняют смысл:
\[
\frac{x+1}{-3}=-\frac{x+1}{3}=\frac{-(x+1)}{3}.
\]
\beware{Правило:} знак принадлежит всему коэффициенту, например $-4x=(-4)\cdot x$.

\sect{8. Правило «ВУПС»: скобки защищают выражение}
На занятии использовали памятку:
\[
\boxed{\text{Вычитаемое, Умножаемое, Подставляемое --- в Скобки.}}
\]
Особенно важно заключать в скобки выражение, состоящее из нескольких слагаемых.

\textbf{Вычитаемое.} Разность выражений $2x+3$ и $4x-7$:
\[
(2x+3)-(4x-7)=2x+3-4x+7=\boxed{-2x+10}.
\]
Внешний минус меняет \emph{оба} знака во второй скобке.

\textbf{Умножаемое.} Их произведение записывается так:
\[
\boxed{(2x+3)(4x-7)}.
\]
Скобки показывают, что каждый множитель --- целое выражение.

\textbf{Подставляемое.} При $x=2y-1$:
\[
2x+3=2(2y-1)+3=4y-2+3=\boxed{4y+1}.
\]

\sect{9. Распределительный закон}
\[
\boxed{k(a+b)=ka+kb},\qquad \boxed{k(a-b)=ka-kb}.
\]
\textbf{Пример из уравнения:}
\[
10(4x+1)=10\cdot4x+10\cdot1=40x+10.
\]
Число $10$ умножается на два \emph{слагаемых}: $4x$ и $1$. Слагаемое $4x$ воспринимается целиком.

\clearpage
\section*{Тренировочный блок. Путь Б}
\textcolor{muted}{27 октября -- 1 ноября 2026 г. · 6 дней · 10 заданий ежедневно.}

Решайте задания в указанном порядке. Для дробных уравнений сначала указывайте допустимые значения знаменателя. В заданиях 8--10 каждого дня обязательно правильно расставляйте скобки по правилу «ВУПС». В ответах к таким заданиям показана и правильная запись со скобками.
\subsection*{День 1. Вторник, 27 октября 2026 г.}
\textcolor{muted}{Сначала выполните задания самостоятельно, затем сверьтесь с таблицей ответов.}
\begin{enumerate}[label=\arabic*.,leftmargin=1.8em,itemsep=11pt,topsep=7pt]
\item Найдите вычитаемое $b$: $23-b=14$.
\item Найдите $x$ из пропорции $\displaystyle\frac35=\frac{x}{20}$.
\item Из формулы $\displaystyle F=G\frac{m_1m_2}{R^2}$ выразите $m_2$.
\item Решите уравнение $5x+7=2x+19$.
\item Решите уравнение $\displaystyle\frac{3x+2}{4}=5$.
\item Решите уравнение $\displaystyle\frac{6}{2x+1}=3$.
\item Решите уравнение $\displaystyle\frac{8}{-3x+5}=4$.
\item \textbf{ВУПС: вычитаемое.} Запишите со скобками и упростите разность выражений $3x+5$ и $2x-7$.
\item \textbf{ВУПС: умножаемое.} Запишите со скобками произведение выражений $x+2$ и $3x-1$ (раскрывать скобки не нужно).
\item \textbf{ВУПС: подставляемое.} Подставьте $x=2y-3$ в $4x+1$ и упростите.
\end{enumerate}
\vfill\small\textcolor{muted}{Самопроверка: знаки, скобки, ограничения на знаменатели, подстановка ответа.}\normalsize
\clearpage
\subsection*{День 2. Среда, 28 октября 2026 г.}
\textcolor{muted}{Сначала выполните задания самостоятельно, затем сверьтесь с таблицей ответов.}
\begin{enumerate}[label=\arabic*.,leftmargin=1.8em,itemsep=11pt,topsep=7pt]
\item Найдите множитель $b$: $8b=56$.
\item Найдите $x$ из пропорции $\displaystyle\frac{x}{9}=\frac43$.
\item Из формулы $\displaystyle F=G\frac{m_1m_2}{R^2}$ выразите $R^2$.
\item Решите уравнение $9x-4=5x+16$.
\item Решите уравнение $\displaystyle\frac{6x-7}{5}=7$.
\item Решите уравнение $\displaystyle\frac{12}{3x-2}=4$.
\item Решите уравнение $\displaystyle\frac{10}{-2x+9}=2$.
\item \textbf{ВУПС: вычитаемое.} Запишите со скобками и упростите разность выражений $7x-2$ и $3x+8$.
\item \textbf{ВУПС: умножаемое.} Запишите произведение числа $5$ и выражения $2a-3$ со скобками, затем раскройте их.
\item \textbf{ВУПС: подставляемое.} Подставьте $a=3t+2$ в $2a-5$ и упростите.
\end{enumerate}
\vfill\small\textcolor{muted}{Самопроверка: знаки, скобки, ограничения на знаменатели, подстановка ответа.}\normalsize
\clearpage
\subsection*{День 3. Четверг, 29 октября 2026 г.}
\textcolor{muted}{Сначала выполните задания самостоятельно, затем сверьтесь с таблицей ответов.}
\begin{enumerate}[label=\arabic*.,leftmargin=1.8em,itemsep=11pt,topsep=7pt]
\item Найдите делимое $a$: $\displaystyle\frac a7=9$.
\item Найдите $x$ из пропорции $\displaystyle\frac5x=\frac{10}{18}$.
\item Из формулы $\displaystyle 2R=\frac{Gm_1m_2}{F}$ выразите $R$.
\item Решите уравнение $12-3x=5x-4$.
\item Решите уравнение $\displaystyle\frac{7x+3}{6}=5$.
\item Решите уравнение $\displaystyle\frac9{2x-5}=3$.
\item Решите уравнение $\displaystyle\frac{15}{-5x+10}=3$.
\item \textbf{ВУПС: вычитаемое.} Запишите со скобками и упростите разность выражений $4x-9$ и $-2x+5$.
\item \textbf{ВУПС: умножаемое.} Запишите со скобками произведение $-3$ и $2p-7$, затем раскройте скобки.
\item \textbf{ВУПС: подставляемое.} Подставьте $x=5-2y$ в $3x-4$ и упростите.
\end{enumerate}
\vfill\small\textcolor{muted}{Самопроверка: знаки, скобки, ограничения на знаменатели, подстановка ответа.}\normalsize
\clearpage
\subsection*{День 4. Пятница, 30 октября 2026 г.}
\textcolor{muted}{Сначала выполните задания самостоятельно, затем сверьтесь с таблицей ответов.}
\begin{enumerate}[label=\arabic*.,leftmargin=1.8em,itemsep=11pt,topsep=7pt]
\item Найдите делитель $b$: $\displaystyle\frac{42}{b}=6$.
\item Найдите $x$ из пропорции $\displaystyle\frac38=\frac9x$.
\item Из формулы $\displaystyle F=G\frac{m_1m_2}{R^2}$ выразите $G$.
\item Решите уравнение $7(x-2)=3x+10$.
\item Решите уравнение $\displaystyle\frac{5x-3}{7}=6$.
\item Решите уравнение $\displaystyle\frac8{5x+3}=2$.
\item Решите уравнение $\displaystyle\frac{12}{-4x-2}=-3$.
\item \textbf{ВУПС: вычитаемое.} Запишите со скобками и упростите разность выражений $5x+2$ и $x-9$.
\item \textbf{ВУПС: умножаемое.} Запишите со скобками произведение $-4$ и $3a-2$, затем раскройте скобки.
\item \textbf{ВУПС: подставляемое.} Подставьте $x=3z-4$ в $5-2x$ и упростите.
\end{enumerate}
\vfill\small\textcolor{muted}{Самопроверка: знаки, скобки, ограничения на знаменатели, подстановка ответа.}\normalsize
\clearpage
\subsection*{День 5. Суббота, 31 октября 2026 г.}
\textcolor{muted}{Сначала выполните задания самостоятельно, затем сверьтесь с таблицей ответов.}
\begin{enumerate}[label=\arabic*.,leftmargin=1.8em,itemsep=11pt,topsep=7pt]
\item Найдите вычитаемое $b$: $50-b=17$.
\item Найдите $x$ из пропорции $\displaystyle\frac{x-1}{6}=\frac54$.
\item Из формулы $\displaystyle F=G\frac{m_1m_2}{R^2}$ выразите $R$ при положительных величинах.
\item Решите уравнение $-4x+9=3x-12$.
\item Решите уравнение $\displaystyle\frac{2x-5}{3}=7$.
\item Решите уравнение $\displaystyle\frac7{3x-4}=1$.
\item Решите уравнение $\displaystyle\frac6{-2x+1}=-3$.
\item \textbf{ВУПС: вычитаемое.} Запишите со скобками и упростите разность выражений $2x-11$ и $5-3x$.
\item \textbf{ВУПС: умножаемое.} Запишите со скобками произведение $7$ и $2t-5$, затем раскройте скобки.
\item \textbf{ВУПС: подставляемое.} Подставьте $t=1-4k$ в $6-3t$ и упростите.
\end{enumerate}
\vfill\small\textcolor{muted}{Самопроверка: знаки, скобки, ограничения на знаменатели, подстановка ответа.}\normalsize
\clearpage
\subsection*{День 6. Воскресенье, 1 ноября 2026 г.}
\textcolor{muted}{Сначала выполните задания самостоятельно, затем сверьтесь с таблицей ответов.}
\begin{enumerate}[label=\arabic*.,leftmargin=1.8em,itemsep=11pt,topsep=7pt]
\item Найдите делитель $b$: $\displaystyle\frac{48}{b}=12$.
\item Найдите $x$ из пропорции $\displaystyle\frac{2x+1}{5}=\frac73$.
\item Из формулы $\displaystyle F=G\frac{m_1m_2}{R^2}$ выразите произведение $m_1m_2$.
\item Решите уравнение $4(2x-3)-5=3(x+1)$.
\item Решите уравнение $\displaystyle\frac{5x+4}{6}=7$.
\item Решите уравнение $\displaystyle\frac{15}{4x-1}=5$.
\item Решите уравнение $\displaystyle\frac9{-3x+6}=-3$.
\item \textbf{ВУПС: вычитаемое.} Запишите со скобками и упростите разность выражений $-x+6$ и $2x-9$.
\item \textbf{ВУПС: умножаемое.} Запишите со скобками произведение $-2$ и $4a-9$, затем раскройте скобки.
\item \textbf{ВУПС: подставляемое.} Подставьте $x=2y-5$ в $3(x-1)-2x$ и упростите.
\end{enumerate}
\vfill\small\textcolor{muted}{Самопроверка: знаки, скобки, ограничения на знаменатели, подстановка ответа.}\normalsize
\clearpage
\section*{Ответы к тренировочному блоку}
\textcolor{muted}{Для задач на ВУПС приводится запись со скобками и итог.}
\subsection*{День 1. Вторник, 27 октября 2026 г.}
\begingroup
\small
\renewcommand{\arraystretch}{1.45}
\begin{tabularx}{\linewidth}{@{}>{\centering\arraybackslash}p{10mm}X@{}}
\toprule
\textbf{№} & \textbf{Ответ / корректная запись} \\
\midrule
1 & $b=9$ \\
2 & $x=12$ \\
3 & $\displaystyle m_2=\frac{FR^2}{Gm_1}$ \\
4 & $x=4$ \\
5 & $x=6$ \\
6 & $\displaystyle x=\frac12$ \\
7 & $x=1$ \\
8 & $(3x+5)-(2x-7)=x+12$ \\
9 & $(x+2)(3x-1)$ \\
10 & $4(2y-3)+1=8y-11$ \\
\bottomrule
\end{tabularx}
\endgroup
\vspace{8mm}
\subsection*{День 2. Среда, 28 октября 2026 г.}
\begingroup
\small
\renewcommand{\arraystretch}{1.45}
\begin{tabularx}{\linewidth}{@{}>{\centering\arraybackslash}p{10mm}X@{}}
\toprule
\textbf{№} & \textbf{Ответ / корректная запись} \\
\midrule
1 & $b=7$ \\
2 & $x=12$ \\
3 & $\displaystyle R^2=\frac{Gm_1m_2}{F}$ \\
4 & $x=5$ \\
5 & $x=7$ \\
6 & $\displaystyle x=\frac53$ \\
7 & $x=2$ \\
8 & $(7x-2)-(3x+8)=4x-10$ \\
9 & $5(2a-3)=10a-15$ \\
10 & $2(3t+2)-5=6t-1$ \\
\bottomrule
\end{tabularx}
\endgroup
\clearpage
\subsection*{День 3. Четверг, 29 октября 2026 г.}
\begingroup
\small
\renewcommand{\arraystretch}{1.45}
\begin{tabularx}{\linewidth}{@{}>{\centering\arraybackslash}p{10mm}X@{}}
\toprule
\textbf{№} & \textbf{Ответ / корректная запись} \\
\midrule
1 & $a=63$ \\
2 & $x=9$ \\
3 & $\displaystyle R=\frac{Gm_1m_2}{2F}$ \\
4 & $x=2$ \\
5 & $\displaystyle x=\frac{27}{7}$ \\
6 & $x=4$ \\
7 & $x=1$ \\
8 & $(4x-9)-(-2x+5)=6x-14$ \\
9 & $-3(2p-7)=-6p+21$ \\
10 & $3(5-2y)-4=11-6y$ \\
\bottomrule
\end{tabularx}
\endgroup
\vspace{8mm}
\subsection*{День 4. Пятница, 30 октября 2026 г.}
\begingroup
\small
\renewcommand{\arraystretch}{1.45}
\begin{tabularx}{\linewidth}{@{}>{\centering\arraybackslash}p{10mm}X@{}}
\toprule
\textbf{№} & \textbf{Ответ / корректная запись} \\
\midrule
1 & $b=7$ \\
2 & $x=24$ \\
3 & $\displaystyle G=\frac{FR^2}{m_1m_2}$ \\
4 & $x=6$ \\
5 & $x=9$ \\
6 & $\displaystyle x=\frac15$ \\
7 & $\displaystyle x=\frac12$ \\
8 & $(5x+2)-(x-9)=4x+11$ \\
9 & $-4(3a-2)=-12a+8$ \\
10 & $5-2(3z-4)=13-6z$ \\
\bottomrule
\end{tabularx}
\endgroup
\clearpage
\subsection*{День 5. Суббота, 31 октября 2026 г.}
\begingroup
\small
\renewcommand{\arraystretch}{1.45}
\begin{tabularx}{\linewidth}{@{}>{\centering\arraybackslash}p{10mm}X@{}}
\toprule
\textbf{№} & \textbf{Ответ / корректная запись} \\
\midrule
1 & $b=33$ \\
2 & $\displaystyle x=\frac{17}{2}$ \\
3 & $\displaystyle R=\sqrt{\frac{Gm_1m_2}{F}}$ \\
4 & $x=3$ \\
5 & $x=13$ \\
6 & $\displaystyle x=\frac{11}{3}$ \\
7 & $\displaystyle x=\frac32$ \\
8 & $(2x-11)-(5-3x)=5x-16$ \\
9 & $7(2t-5)=14t-35$ \\
10 & $6-3(1-4k)=3+12k$ \\
\bottomrule
\end{tabularx}
\endgroup
\vspace{8mm}
\subsection*{День 6. Воскресенье, 1 ноября 2026 г.}
\begingroup
\small
\renewcommand{\arraystretch}{1.45}
\begin{tabularx}{\linewidth}{@{}>{\centering\arraybackslash}p{10mm}X@{}}
\toprule
\textbf{№} & \textbf{Ответ / корректная запись} \\
\midrule
1 & $b=4$ \\
2 & $\displaystyle x=\frac{16}{3}$ \\
3 & $\displaystyle m_1m_2=\frac{FR^2}{G}$ \\
4 & $x=4$ \\
5 & $\displaystyle x=\frac{38}{5}$ \\
6 & $x=1$ \\
7 & $x=3$ \\
8 & $(-x+6)-(2x-9)=-3x+15$ \\
9 & $-2(4a-9)=-8a+18$ \\
10 & $3((2y-5)-1)-2(2y-5)=2y-8$ \\
\bottomrule
\end{tabularx}
\endgroup
\vfill
\begin{center}
\textcolor{darkaccent}{\rule{.30\linewidth}{.5pt}}\\[4mm]
\small Лёвин Артём Александрович — эксклюзивно для Володи Хачатуряна
\end{center}
\end{document}
